\title{: Advancing the Boundaries of Mathematical Reasoning in Open Language Models}

\begin{abstract}
Mathematical reasoning represents a considerable challenge for language models because of its intricate and structured characteristics. In this work, we present DeepSeekMath~7B, which extends the pre-training of DeepSeek-Coder-Base-v1.5 7B by incorporating 120B math-focused tokens derived from Common Crawl, alongside natural language and code datasets. DeepSeekMath~7B attains a notable 51.7\% score on the competition-level MATH benchmark without the use of external tools or ensemble voting methods, nearing the performance of Gemini-Ultra and GPT-4. Utilizing self-consistency with 64 samples from DeepSeekMath~7B results in a 60.9\% accuracy on MATH. The enhanced mathematical reasoning ability of DeepSeekMath~is credited to two primary factors: first, leveraging the substantial potential of publicly accessible web data through a carefully designed data selection pipeline; second, introducing Group Relative Policy Optimization (GRPO), a PPO variant that improves mathematical reasoning skills while simultaneously reducing PPO's memory consumption.
\end{abstract}

\section{Introduction}

Large language models (LLMs) have transformed the field of mathematical reasoning within artificial intelligence, driving major progress on benchmarks for quantitative reasoning \citep{MATH} as well as geometric reasoning \citep{trinh2024solving}. Furthermore, these models have demonstrated their utility in aiding humans to tackle challenging mathematical problems \citep{tao}. Nevertheless, state-of-the-art models such as GPT-4 \citep{gpt4} and Gemini-Ultra \citep{gemini} remain proprietary, while currently available open-source models lag significantly behind in terms of capability.

In this work, we present DeepSeekMath, a domain-specific language model that considerably surpasses the mathematical performance of existing open-source models and nears the capabilities of GPT-4 on scholarly benchmarks. To accomplish this, we develop the DeepSeekMath Corpus, a large-scale, high-quality pre-training dataset containing 120B math-related tokens. This dataset is derived from the Common Crawl (CC) by utilizing a fastText-based classifier \citep{joulin2016fasttext}. Initially, the classifier is trained with positive examples from OpenWebMath \citep{openwebmath}, alongside a varied collection of other web pages as negative examples. Afterward, the classifier is used to identify additional positive samples from the CC, which are subsequently refined through human annotation. The classifier is then retrained with this improved dataset to enhance its accuracy. Evaluation results demonstrate that the large-scale corpus maintains high quality, with our base model DeepSeekMath-Base 7B attaining 64.2\% on GSM8K \citep{gsm8k} and 36.2\% on the challenging MATH dataset \citep{MATH}, outperforming Minerva 540B \citep{minerva}. Furthermore, since the DeepSeekMath Corpus is multilingual, we observe gains on Chinese mathematical benchmarks \citep{wei2023cmath,agieval}. We consider our efforts in mathematical data processing to serve as a foundation for the research community, acknowledging that there remains substantial potential for further enhancements.

DeepSeekMath-Base is initialized using DeepSeek-Coder-Base-v1.5 7B \citep{deepseek-coder}, as we find that beginning with a model trained on code is more advantageous than starting from a general LLM. Additionally, we observe that training on mathematical tasks also enhances the model's performance on MMLU \citep{mmlu} and BBH benchmarks \citep{bbh}, suggesting that this training not only boosts the model’s math skills but also strengthens its overall reasoning abilities.

Following pre-training, we perform mathematical instruction tuning on DeepSeekMath-Base utilizing chain-of-thought \citep{cot}, program-of-thought \citep{pot,pal}, and tool-integrated reasoning \citep{tora} datasets. The resulting model, DeepSeekMath-Instruct 7B, outperforms all other 7B models and achieves performance comparable to 70B open-source instruction-tuned models.

Moreover, we propose Group Relative Policy Optimization (GRPO), a variant of the Proximal Policy Optimization (PPO) \citep{schulman2017proximal} reinforcement learning algorithm. GRPO eliminates the use of a critic model, instead estimating the baseline from group scores, which substantially lowers the computational resources required for training. By employing only a subset of English instruction tuning data, GRPO achieves significant improvements over the strong DeepSeekMath-Instruct model, including enhancements on both in-domain tasks (GSM8K: 82.9\% $\rightarrow$ 88.2\%, MATH: 46.8\% $\rightarrow$ 51.7\%) and out-of-domain mathematical challenges (e.g., CMATH: 84.6\% $\rightarrow$ 88.8\%) during reinforcement learning. We also present a unified framework to interpret various methods, such as Rejection Sampling Fine-Tuning (RFT) \citep{yuan2023scaling}, Direct Preference Optimization (DPO) \citep{dpo}, PPO, and GRPO. Within this framework, we identify all these approaches as either direct or simplified reinforcement learning techniques. We conduct comprehensive experiments, including comparisons of online versus offline training, outcome versus process supervision, and single-turn versus iterative RL, to thoroughly explore the fundamental components of this framework. Finally, we discuss why our RL approach enhances the performance of instruction-tuned models and summarize future directions for developing more effective RL methods based on this unified framework.

\subsection{Contributions}
Our work presents scalable math pre-training in addition to investigating and analyzing reinforcement learning.

\textbf{Math Pre-Training at Scale}

\begin{itemize}[topsep=0pt]
    \item Our study demonstrates strong evidence that publicly available Common Crawl data holds significant value for mathematical applications. Through a carefully designed data selection process, we create the DeepSeekMath Corpus, a high-quality dataset comprising 120B tokens extracted from web pages curated for mathematical content. This dataset is nearly 7 times larger than the math web pages used by Minerva \citep{minerva} and 9 times larger than the recently introduced OpenWebMath \citep{openwebmath}.

    \item Our pre-trained base model, DeepSeekMath-Base 7B, attains performance on par with Minerva 540B \citep{minerva}, suggesting that model size alone is not the sole determinant of mathematical reasoning capabilities. A smaller model trained on high-quality data can also achieve strong results.

    \item We present insights from our math training experiments. Pre-training on code before math training enhances models’ performance on mathematical problem-solving tasks both with and without the use of tools. This provides a partial resolution to the longstanding question: \textit{does code training enhance reasoning skills?} Our findings indicate it does, at least within the context of mathematical reasoning.

    \item While training on arXiv papers is a common practice, especially in many math-focused studies, it does not yield significant improvements across the mathematical benchmarks used in this research.
\end{itemize}

\textbf{Exploration and Analysis of Reinforcement Learning}

\begin{itemize}[topsep=0pt]
    \item We propose Group Relative Policy Optimization (GRPO), a reinforcement learning approach that is both efficient and effective. GRPO eliminates the need for a critic model by estimating the baseline using group scores, which greatly reduces the computational resources required compared to Proximal Policy Optimization (PPO).
    \item We show that GRPO notably improves the performance of our instruction-tuned model, DeepSeekMath-Instruct, using only instruction-tuning data. Additionally, we observe improvements in out-of-domain generalization throughout the reinforcement learning process.
    \item We introduce a unified framework to interpret various techniques, including RFT, DPO, PPO, and GRPO. We conduct extensive experiments—such as comparing online versus offline training, outcome versus process supervision, and single-turn versus iterative reinforcement learning—to thoroughly analyze the core aspects of this framework.
    \item Leveraging our unified framework, we investigate the underlying causes of reinforcement learning’s effectiveness and outline several promising directions to enhance reinforcement learning for large language models (LLMs).
\end{itemize}

\subsection{Summary of Evaluations and Metrics}

\begin{itemize}[topsep=0pt]
    \item \textbf{English and Chinese Mathematical Reasoning}: We perform comprehensive evaluations of our models on both English and Chinese benchmarks, spanning mathematical problems from elementary to university levels. The English datasets include GSM8K \citep{gsm8k}, MATH \citep{MATH}, SAT \citep{llemma}, OCW Courses \citep{minerva}, and MMLU-STEM \citep{mmlu}. The Chinese datasets comprise MGSM-zh \citep{mgsm}, CMATH \citep{wei2023cmath}, Gaokao-MathCloze \citep{agieval}, and Gaokao-MathQA \citep{agieval}. We assess models’ capability to generate complete textual solutions without external tools as well as their problem-solving skills using Python.
    
    On English benchmarks, DeepSeekMath-Base performs on par with the closed-source Minerva 540B \citep{minerva} and outperforms all open-source base models (e.g., Mistral 7B \citep{mistral} and Llemma-34B \citep{llemma}), regardless of prior math-focused pre-training, often by a considerable margin. Remarkably, DeepSeekMath-Base shows even stronger results on Chinese benchmarks, likely because, unlike prior works \citep{minerva,llemma}, we do not restrict math pre-training data to English only and incorporate high-quality non-English data. With mathematical instruction tuning and reinforcement learning, the resulting DeepSeekMath-Instruct and DeepSeekMath-RL achieve impressive performance, surpassing 50\% accuracy on the competition-level MATH dataset—a first for open-source models.
\end{itemize}

\begin{itemize}
    \item \textbf{Formal Mathematics}: We assess DeepSeekMath-Base on the informal-to-formal theorem proving task introduced by \citep{dsp_proof} using the miniF2F dataset \citep{minif2f} with Isabelle \citep{isabelle} as the chosen proof assistant. DeepSeekMath-Base exhibits strong few-shot capabilities for autoformalization.
    \item \textbf{Natural Language Understanding, Reasoning, and Code}: To create a detailed assessment of the model's overall comprehension, reasoning, and programming abilities, we evaluate DeepSeekMath-Base on the Massive Multitask Language Understanding (MMLU) benchmark \citep{mmlu}, which includes 57 multiple-choice tasks spanning a wide range of subjects; BIG-Bench Hard (BBH) \citep{bbh}, consisting of 23 difficult tasks primarily requiring multi-step reasoning; as well as HumanEval \citep{codex} and MBPP \citep{mbpp}, commonly used benchmarks for evaluating code language models. Pre-training on mathematical data improves both language comprehension and reasoning outcomes.
\end{itemize}

\section{Math Pre-Training}

\subsection{Data Collection and Decontamination} This section details the construction process of the DeepSeekMath Corpus from Common Crawl data. As illustrated in Figure \ref{fig:data_collect}, we propose an iterative methodology to systematically accumulate a large-scale mathematical corpus from Common Crawl, beginning with a seed dataset (for example, a smaller yet high-quality math-focused dataset). Notably, this method can also be adapted for other domains like programming.

Initially, we select OpenWebMath \citep{openwebmath}, a high-quality set of mathematical web texts, as our seed corpus. From this dataset, we train a fastText model \citep{joulin2016fasttext} to identify additional web pages resembling OpenWebMath’s mathematical content. Specifically, we randomly pick 500,000 samples from the seed corpus as positive examples and 500,000 Common Crawl web pages as negatives for training. Using an open-source library\footnote{\url{https://fasttext.cc}}, the model is trained with a vector size of 256, learning rate 0.1, maximum word n-gram length of 3, minimum word frequency of 3, and 3 training epochs. To reduce the original Common Crawl size, we apply URL-based deduplication and near-deduplication techniques, yielding 40 billion HTML pages. We then use the fastText model to retrieve mathematical web pages from this deduplicated set. To eliminate lower-quality mathematical content, pages are ranked by their fastText prediction scores, and only the highest-ranking ones are retained. The amount of retained data is evaluated via pre-training experiments using the top 40B, 80B, 120B, and 160B tokens. In the first cycle, we opt to keep the top 40B tokens.

Following the initial round of data collection, many mathematical web pages remain uncollected, primarily because the fastText model was trained on a set of positive examples that lacked sufficient diversity. To address this, we identify more mathematical web sources to expand the seed corpus, enabling us to fine-tune the fastText model further. Concretely, we first partition the entire Common Crawl into separate domains, where a domain is defined as web pages sharing the same base URL. For each domain, we compute the proportion of web pages collected in the first iteration. Domains with over 10\% of their pages collected are classified as math-related (e.g., \url{mathoverflow.net}). 

Next, we manually label URLs within these identified domains that contain mathematical content (e.g., \url{mathoverflow.net/questions}). Web pages linked to these URLs, but not yet collected, are then added to the seed corpus. This strategy allows us to accumulate more positive examples, facilitating the training of an enhanced fastText model capable of retrieving more mathematical data in the next iteration. After four iterations, we accumulate 35.5M mathematical web pages comprising 120B tokens. Notably, during the fourth iteration, we observe that nearly 98\% of the data had already been captured by the third iteration, leading us to halt the data collection process.

To prevent contamination of benchmarks, we adhere to \cite{deepseek-coder} by excluding web pages containing questions or answers from English mathematical benchmarks such as GSM8K \citep{gsm8k} and MATH \citep{MATH}, as well as Chinese benchmarks like CMATH \citep{wei2023cmath} and AGIEval \citep{agieval}. The filtering procedure involves removing any text segment containing a 10-gram sequence that exactly matches any substring from these evaluation benchmarks. For benchmark texts shorter than 10 grams but at least 3 grams long, we apply exact matching to exclude contaminated pages from our math training corpus.

\subsection{Validating the Quality of the DeepSeekMath~Corpus}

We conduct pre-training experiments to compare the DeepSeekMath~Corpus against recently released math-training corpora:

\begin{itemize}[topsep=0pt]
    \item \textbf{MathPile} \citep{mathpile}: a multi-source corpus containing 8.9B tokens aggregated from textbooks, Wikipedia, ProofWiki, CommonCrawl, StackExchange, and arXiv, with over 85\% of the data derived from arXiv;
    \item \textbf{OpenWebMath} \citep{openwebmath}: filtered CommonCrawl data focusing on mathematical content, totaling 13.6B tokens;
    \item \textbf{Proof-Pile-2} \citep{llemma}: a mathematical corpus comprising OpenWebMath, AlgebraicStack (10.3B tokens of mathematical code), and arXiv papers (28.0B tokens). When conducting experiments on Proof-Pile-2, we follow \cite{llemma} by using an arXiv:Web:Code data ratio of 2:4:1.
\end{itemize}

\subsubsection{Training Configuration}
\label{sec:quality-policy}
We conduct math training on a general pre-trained language model with 1.3 billion parameters, built on the same architecture as the DeepSeek LLMs \citep{deepseek-llm}, referred to as DeepSeek-LLM 1.3B. Each mathematical corpus is used to train a separate model for 150 billion tokens. All experiments utilize the efficient and lightweight HAI-LLM \citep{haillm} training framework. Adhering to the DeepSeek LLMs training protocol, we employ the AdamW optimizer \citep{adamW} with $\beta_1=0.9$, $\beta_2=0.95$, and $\mathrm{weight\_decay}=0.1$, along with a multi-phase learning rate schedule: the learning rate peaks after 2,000 warm-up steps, then decays to 31.6\% of the peak at 80\% progress of training, and further drops to 10.0\% of the peak at 90\% of the training duration. The maximum learning rate is set to 5.3e-4, and the batch size is fixed at 4 million tokens with a 4,000 token context window.

\subsubsection{Evaluation Outcomes}

\textbf{The DeepSeekMath~Corpus exhibits superior quality, contains multilingual mathematical data, and is the largest corpus available.}

\begin{itemize}[topsep=0pt]
    \item \textbf{Superior Quality}:
    We assess downstream performance on eight mathematical benchmarks using few-shot chain-of-thought prompting \cite{cot}. As presented in Table \ref{tab:corpora_comparison}, the model trained on the DeepSeekMath~Corpus consistently outperforms others. Figure \ref{fig:corpora_comparisons} illustrates that after 50 billion tokens (equivalent to one full epoch of Proof-Pile-2), the DeepSeekMath-trained model surpasses the one trained on Proof-Pile-2, suggesting that the average quality of the DeepSeekMath Corpus is higher.
    \item \textbf{Multilingual}:
    The DeepSeekMath~Corpus incorporates content in various languages, with English and Chinese being the most prevalent. As shown in Table \ref{tab:corpora_comparison}, training on this corpus improves mathematical reasoning performance in both English and Chinese. By contrast, existing mathematical corpora, which primarily focus on English, demonstrate limited or even detrimental effects on Chinese mathematical reasoning tasks.
    \item \textbf{Large Scale}:
    The DeepSeekMath~Corpus is significantly larger than other mathematical corpora currently available. Figure \ref{fig:corpora_comparisons} depicts that DeepSeek-LLM 1.3B, trained on the DeepSeekMath~Corpus, exhibits a steeper learning trajectory and more sustained performance gains. Conversely, the smaller baseline corpora have been cycled through multiple training rounds, leading the model’s performance to quickly plateau.
\end{itemize}

\subsection{Training and Assessing DeepSeekMath-Base 7B}

Here, we present DeepSeekMath-Base 7B, a foundational model with notable reasoning capabilities, particularly in the domain of mathematics. The model is initialized from DeepSeek-Coder-Base-v1.5 7B \citep{deepseek-coder} and trained on a corpus of 500B tokens. The data composition is as follows: 56\% originates from the DeepSeekMath Corpus, 4\% from AlgebraicStack, 10\% from arXiv, 20\% consists of GitHub code, and the remaining 10\% includes natural language content from Common Crawl in both English and Chinese. We primarily follow the training protocol detailed in Section \ref{sec:quality-policy}, with adjustments to a maximum learning rate of 4.2e-4 and a batch size of 10M tokens.

We perform an extensive evaluation of DeepSeekMath-Base 7B’s mathematical prowess, emphasizing its capability to generate autonomous mathematical solutions without external aids, tackle problems utilizing tools, and carry out formal theorem proving. Besides mathematical tasks, we also offer a broader evaluation of the base model’s general abilities, such as natural language comprehension, logical reasoning, and coding proficiency.

\paragraph{Stepwise Reasoning for Mathematical Problem Solving}

We assess DeepSeekMath-Base’s problem-solving effectiveness using few-shot chain-of-thought prompting \citep{cot} over eight benchmarks in both English and Chinese. These datasets include quantitative reasoning tests (e.g., GSM8K \citep{gsm8k}, MATH \citep{MATH}, and CMATH \citep{wei2023cmath}) as well as multiple-choice challenges (e.g., MMLU-STEM \citep{mmlu} and Gaokao-MathQA \citep{agieval}), spanning a wide spectrum of mathematical topics from elementary to collegiate difficulty levels.

As indicated in Table \ref{tab:base_math_cot}, DeepSeekMath-Base 7B achieves the highest performance across all eight benchmarks among the open-source base models, which includes the widely adopted general model Mistral 7B \citep{mistral} and the recently introduced Llemma 34B \citep{llemma} that was trained on Proof-Pile-2 \citep{llemma} for mathematical tasks. Remarkably, on the challenging competition-level MATH dataset, DeepSeekMath-Base surpasses existing open-source base models by more than 10\% absolute and even outperforms Minerva 540B \citep{minerva}, a closed-source base model that is 77 times larger, based on PaLM \citep{palm} and further trained on specialized mathematical literature.

\paragraph{Mathematical Problem Solving with Tool Use} We assess program-assisted mathematical reasoning on GSM8K and MATH datasets utilizing few-shot program-of-thought prompting \citep{pot,pal}. Models are prompted to address each problem by generating a Python program that may use libraries such as \textit{math} and \textit{sympy} for complex calculations. The outcome of the program execution is considered as the answer. As presented in Table \ref{tab:base_math_pot_proof}, DeepSeekMath-Base 7B exceeds the performance of the previous state-of-the-art model Llemma 34B.

\paragraph{Formal Mathematics}

Formal proof automation plays an important role in verifying the correctness and dependability of mathematical proofs while improving efficiency, and has received growing focus in recent times. We test DeepSeekMath-Base 7B on the informal-to-formal proving task from \citep{dsp_proof}, which involves producing a formal proof based on an informal statement, its formal equivalent, and the informal proof itself. The evaluation is conducted on miniF2F \citep{minif2f}, a benchmark targeting formal Olympiad-level mathematics, where a formal proof in Isabelle is generated for each problem using few-shot prompting. Following the approach in \cite{dsp_proof}, models generate proof outlines, which are then completed using the automated prover Sledgehammer \citep{sledgehammer}. As demonstrated in Table \ref{tab:base_math_pot_proof}, DeepSeekMath-Base 7B exhibits robust performance in automatic proof formalization.

\paragraph{Natural Language Understanding, Reasoning, and Code}

We examine model capabilities in natural language understanding on MMLU \citep{mmlu}, reasoning on BBH \citep{bbh}, and coding on HumanEval \citep{codex} and MBPP \citep{mbpp}. As shown in Table \ref{tab:understanding_reasoning_code}, DeepSeekMath-Base 7B shows notable improvements over its predecessor, DeepSeek-Coder-Base-v1.5 \citep{deepseek-coder}, in MMLU and BBH performance, highlighting the beneficial effect of mathematical training on language understanding and reasoning skills. Moreover, by incorporating code tokens during continual training, DeepSeekMath-Base 7B successfully preserves the coding benchmark performance of DeepSeek-Coder-Base-v1.5. Overall, DeepSeekMath-Base 7B significantly outperforms the general model Mistral 7B \citep{mistral} across the three reasoning and coding tasks.

\section{Supervised Fine-Tuning}

\subsection{SFT Data Curation}

We create a mathematical instruction-tuning dataset that encompasses problems in English and Chinese from various branches of mathematics and with differing levels of difficulty: each problem is paired with solutions presented in chain-of-thought (CoT) \citep{cot}, program-of-thought (PoT) \citep{pot,pal}, and tool-integrated reasoning formats \citep{tora}. The dataset consists of a total of 776K training examples.

\begin{itemize}[topsep=0pt]
    \item \textbf{English mathematical datasets}: We annotate GSM8K and MATH problems with tool-integrated solutions and include a portion of MathInstruct \citep{MathInstruct}, together with the training data from Lila-OOD \citep{lila}, where problems are solved using either CoT or PoT methods. Our English dataset covers a broad range of mathematical areas such as algebra, probability, number theory, calculus, and geometry.
    \item \textbf{Chinese mathematical datasets}: We gather Chinese K-12 math problems covering 76 subtopics, including linear equations, with solutions provided both in CoT and tool-integrated reasoning formats.
\end{itemize}

\subsection{Training and Evaluating DeepSeekMath-Instruct 7B}

Here, we present DeepSeekMath-Instruct 7B, which is fine-tuned on mathematical instructions starting from DeepSeekMath-Base. Training examples are concatenated randomly up to a maximum context length of 4K tokens. The model is trained for 500 steps using a batch size of 256 and a fixed learning rate of 5e-5.

We assess the mathematical capabilities of the models both without and with tool usage, employing four quantitative reasoning benchmarks in English and Chinese. Our model is compared against the top-performing models available at the time:

\begin{itemize}[topsep=0pt]
    \item \textbf{Closed-source models} include: (1) the GPT series, where GPT-4 \citep{gpt4} and GPT-4 Code Interpreter~\footnote{\url{https://openai.com/blog/chatgpt-plugins\#\#code-interpreter}} represent the most advanced versions, (2) Gemini Ultra and Pro \citep{gemini}, (3) Inflection-2 \citep{inflection-2}, (4) Grok-1~\footnote{\url{https://x.ai/model-card}}, along with models recently introduced by Chinese companies such as (5) Baichuan-3~\footnote{\url{https://www.baichuan-ai.com}}, and (6) the latest GLM-4~\footnote{\url{https://open.bigmodel.cn/dev/api\#glm-4}} from the GLM family \citep{glm}.

These models serve general purposes, with most having undergone multiple alignment steps.  
\item \textbf{Open-source models} encompass: general-purpose models such as (1) DeepSeek-LLM-Chat 67B \citep{deepseek-llm}, (2) Qwen 72B \citep{qwen}, (3) SeaLLM-v2 7B \citep{seallm}, and (4) ChatGLM3 6B \citep{chatglm3}, along with models enhanced for mathematical capabilities, including (5) InternLM2-Math 20B~\footnote{\url{https://github.com/InternLM/InternLM-Math}}, which extends InternLM2 with math-focused training followed by instruction tuning, (6) Math-Shepherd-Mistral 7B which applies PPO training \citep{schulman2017proximal} to Mistral 7B \citep{mistral} using a process-supervised reward model, (7) the WizardMath series \citep{wizardmath}, which enhances mathematical reasoning in Mistral 7B and Llama-2 70B \citep{llama2} through evolve-instruct (a variant of instruction tuning utilizing AI-evolved instructions) combined with PPO training using problems mainly drawn from GSM8K and MATH datasets, (8) MetaMath 70B \citep{metamath}, fine-tuned from Llama-2 70B on an augmented GSM8K and MATH dataset, (9) ToRA 34B \cite{tora}, which fine-tunes CodeLlama 34B for tool-integrated mathematical reasoning, and (10) MAmmoTH 70B \citep{MathInstruct}, which instruction-tunes Llama-2 70B on MathInstruct.  
\end{itemize}

As presented in Table \ref{tab:sft_rl_math}, when evaluated without tool usage, DeepSeekMath-Instruct 7B exhibits robust step-by-step reasoning capabilities. Remarkably, on the challenging MATH dataset, our model outperforms all open-source models and most proprietary models (e.g., Inflection-2 and Gemini Pro) by a margin of at least 9\% in absolute accuracy. This holds even against significantly larger models (e.g., Qwen 72B) or those specifically enhanced via math-focused reinforcement learning (e.g., WizardMath-v1.1 7B). While DeepSeekMath-Instruct is comparable to Chinese proprietary models such as GLM-4 and Baichuan-3 on MATH, it still trails behind GPT-4 and Gemini Ultra.

When models are permitted to employ natural language reasoning in combination with program-based tool usage for solving problems, DeepSeekMath-Instruct 7B attains nearly 60\% accuracy on MATH, exceeding all currently available open-source models. Across other benchmarks, our model remains competitive with DeepSeek-LLM-Chat 67B, the previous leading model which is ten times larger.  
\section{Reinforcement Learning}

\subsection{Group Relative Policy Optimization} Reinforcement learning (RL) has demonstrated effectiveness in further enhancing the mathematical reasoning abilities of LLMs following the Supervised Fine-Tuning (SFT) phase \citep{wang2023math,wizardmath}. In this section, we describe our efficient and effective RL method, termed Group Relative Policy Optimization (GRPO).

\subsubsection{From PPO to GRPO} Proximal Policy Optimization (PPO) \citep{schulman2017proximal} is an actor-critic RL method widely applied during the RL fine-tuning of LLMs \citep{ouyang2022training}. Specifically, it improves LLMs by maximizing the surrogate objective:

\begin{equation}
\footnotesize
    \mathcal{J}_{PPO}(\theta) = \mathbb{E}{[q \sim P(Q), o \sim \pi_{\theta_{old}}(O|q)]} \frac{1}{|o|} \sum_{t=1}^{|o|} \min \left[ \frac{\pi_\theta(o_{t} | q, o_{<t})}{\pi_{\theta_{old}}(o_{t} | q, o_{<t})} A_{t}, \text{clip} \left( \frac{\pi_\theta(o_{t} | q, o_{<t})}{\pi_{\theta_{old}}(o_{t} | q, o_{<t})}, 1 - \varepsilon, 1 + \varepsilon \right)  A_{t} \right] ,
\end{equation}

where $\pi_{\theta}$ and $\pi_{\theta_{old}}$ denote the current and previous policy models, respectively, and $q, o$ represent questions and outputs sampled from the question dataset and the old policy $\pi_{\theta_{old}}$, correspondingly. The parameter $\varepsilon$ is a clipping-related hyper-parameter introduced in PPO to stabilize the training process. The term $A_t$ stands for the advantage, calculated using Generalized Advantage Estimation (GAE) \citep{gae}, based on the sequence of rewards $\{r_{\ge t}\}$ and a learned value function $V_{\psi}$. Therefore, in PPO, it is necessary to train a value function alongside the policy model. To prevent excessive optimization towards the reward model, the conventional method involves incorporating a per-token KL penalty from a reference model into the reward at each token \citep{ouyang2022training}, expressed as:

\begin{equation}
    r_{t} = r_\varphi(q, o_{\le t}) - \beta \log\frac{\pi_{\theta}(o_{t}|q, o_{<t})}{\pi_{ref}(o_{t}|q, o_{<t})},
\label{eq:PPO-reward}
\end{equation}

where $r_\varphi$ is the reward model, $\pi_{ref}$ is the reference model, often the initial SFT model, and $\beta$ is the coefficient for the KL penalty.

Since the value function utilized in PPO is generally another model of similar scale to the policy model, it imposes significant memory and computational demands. Moreover, during RL training, the value function acts as a baseline in the advantage computation to reduce variance. However, in the context of LLMs, rewards from the reward model are usually assigned only to the last token, which makes it challenging to train a value function that provides accurate estimates at each token step. To tackle this issue, as depicted in Figure \ref{fig:grpo}, we introduce Group Relative Policy Optimization (GRPO), which eliminates the requirement for a separate value function approximation as in PPO. Instead, GRPO employs the average reward of multiple sampled outputs generated in response to the same question as the baseline. Concretely, for each question $q$, GRPO samples a group of outputs $\{o_1, o_2, \cdots, o_G\}$ from the old policy $\pi_{\theta_{old}}$ and optimizes the policy model by maximizing the following objective:

\begin{equation}
\footnotesize
\begin{split}
    \mathcal{J}_{GRPO}(\theta) &= \mathbb{E}{[q \sim P(Q), \{o_i\}_{i=1}^G \sim \pi_{\theta_{old}}(O|q)]}  \\
    & \frac{1}{G}\sum_{i=1}^G\frac{1}{|o_i|} \sum_{t=1}^{|o_i|} \left\{ \min \left[ \frac{\pi_\theta(o_{i,t} | q, o_{i,<t})}{\pi_{\theta_{old}}(o_{i,t} | q, o_{i,<t})} \hat{A}_{i,t}, \text{clip} \left( \frac{\pi_\theta(o_{i,t} | q, o_{i,<t})}{\pi_{\theta_{old}}(o_{i,t} | q, o_{i,<t})}, 1 - \varepsilon, 1 + \varepsilon \right)  \hat{A}_{i,t} \right] - \beta \mathbb{D}_{KL}\left[\pi_{\theta} || \pi_{ref}\right]\right\} ,
\end{split}
\label{eq:GRPO-obj}
\end{equation}

where $\varepsilon$ and $\beta$ are hyper-parameters, and $\hat{A}_{i,t}$ represents the advantage computed solely from the relative rewards of outputs within each group, which will be elaborated upon in the subsequent sections. This group-relative approach to computing advantages aligns naturally with the comparative characteristics of reward models, as reward models are generally trained on datasets containing comparisons among outputs for the same question. Furthermore, unlike PPO where the KL penalty is incorporated into the reward, GRPO regularizes by directly adding the KL divergence between the trained policy and the reference policy to the loss function, thus simplifying the calculation of $\hat{A}_{i,t}$. Different from the KL penalty term in (\ref{eq:PPO-reward}), we estimate the KL divergence using the following unbiased estimator \citep{kl_approx}:

\begin{equ

\begin{algorithm}[t]
  \small
  \caption{Iterative Group Relative Policy Optimization}
  \textbf{Input} initial policy model $\pi_{\theta_{\text{init}}}$; reward models $r_\varphi$; task prompts $\mathcal{D}$; 
  hyperparameters $\varepsilon$, $\beta$, $\mu$
  \begin{algorithmic}[1]
    \State Initialize policy model $\pi_\theta$ as $\pi_{\theta_{\text{init}}}$
    \For{iteration = 1, \dots, I}
       \State Set reference model $\pi_{ref}$ to current policy $\pi_{\theta}$
      \For{step = 1, \dots, M}
      \State Draw a batch $\mathcal{D}_b$ from $\mathcal{D}$
      \State Save current policy as old policy $\pi_{\theta_{old}} \leftarrow \pi_{\theta}$ 
      \State For each question $q \in \mathcal{D}_b$, sample $G$ outputs $\{o_i\}_{i=1}^G$ from $\pi_{\theta_{old}}(\cdot \mid q)$
      \State Evaluate the outputs $\{o_i\}_{i=1}^G$ with reward model $r_{\varphi}$ to obtain rewards $\{r_i\}_{i=1}^G$
      \State Calculate $\hat{A}_{i,t}$ for each token $t$ in output $o_i$ via group relative advantage estimation
      \For{GRPO iteration = 1, \dots, $\mu$}
        \State Optimize policy model $\pi_{\theta}$ by maximizing the GRPO objective (Equation \ref{eq:GC-GRPO})
      \EndFor
    \EndFor \State Update reward model parameters $r_\varphi$ continuously using a replay buffer
    \EndFor 
  \end{algorithmic}
  \textbf{Output} updated policy $\pi_\theta$
  \label{alg:iter-grpo}
\end{algorithm}

\subsubsection{Outcome Supervision RL with GRPO} Concretely, for each question $q$, a set of outputs $\{o_1, o_2, \ldots, o_G\}$ are drawn from the previous policy $\pi_{\theta_{old}}$. Each output is then evaluated by the reward model to produce a corresponding reward vector $\mathbf{r}=\{r_1, r_2, \ldots, r_G\}$. These rewards are normalized by subtracting the group mean and dividing by the group standard deviation. Outcome supervision assigns the normalized reward to all tokens within each output, i.e., $\hat{A}_{i,t} = \widetilde{r}_i = \frac{r_i - {\rm mean}(\mathbf{r})}{{\rm std}(\mathbf{r})}$, and then updates the policy by maximizing the objective in equation (\ref{eq:GRPO-obj}).

\subsubsection{Process Supervision RL with GRPO} Since outcome supervision only provides a reward at the completion of the entire output, it might be insufficient or inefficient for guiding the policy in complex mathematical reasoning tasks. Inspired by \cite{wang2023math}, we also investigate process supervision, which assigns rewards after each reasoning step. Formally, given a question $q$ and $G$ sampled outputs $\{o_1, o_2, \ldots, o_G\}$, a process reward model scores every reasoning step, producing rewards: $\mathbf{R} = \{\{r_1^{index(1)},\ldots,r_1^{index(K_1)}\}, \ldots, \{r_G^{index(1)},\ldots,r_G^{index(K_G)}\}\}$, where $index(j)$ denotes the token index marking the end of the $j$-th step, and $K_i$ is the number of steps in output $o_i$. These step-wise rewards are normalized across the group: $\widetilde{r}_i^{index(j)} = \frac{r_i^{index(j)} - {\rm mean}(\mathbf{R})}{{\rm std}(\mathbf{R})}$. The advantage for each token is then computed as the sum of the normalized rewards from all subsequent steps, i.e., $\hat{A}_{i,t} = \sum_{index(j) \ge t} \widetilde{r}_i^{index(j)}$, followed by policy optimization via maximizing the objective in equation (\ref{eq:GRPO-obj}).

\subsubsection{Iterative RL with GRPO} As the reinforcement learning training advances, the previous reward model might no longer effectively guide the current policy model. Hence, we investigate an iterative RL approach using GRPO. As outlined in Algorithm \ref{alg:iter-grpo}, iterative GRPO involves generating new training datasets for the reward model from samples produced by the policy model, and continuously updating the existing reward model through a replay mechanism that integrates 10\% of past data. Subsequently, we designate the policy model as the reference model and keep refining the policy model using the updated reward model.

\subsection{Training and Evaluating DeepSeekMath-RL}

We perform RL training based on DeepSeekMath-Instruct 7B. The RL training data consists of chain-of-thought style questions drawn from GSM8K and MATH within the SFT dataset, totaling approximately 144K questions. Other SFT questions are excluded to specifically assess the influence of RL on benchmarks with limited data during the RL phase. The training dataset for reward models is constructed in accordance with \citep{wang2023math}. Our initial reward model is trained starting from DeepSeekMath-Base 7B using a learning rate of 2e-5. For GRPO, the policy model’s learning rate is set to 1e-6, and the KL divergence coefficient is fixed at 0.04. For each question, we sample $64$ outputs. The maximum sequence length is 1024, and the training batch size is set to 1024. The policy model undergoes a single update after each exploration phase. We evaluate DeepSeekMath-RL 7B on benchmarks in the same manner as DeepSeekMath-Instruct 7B. For DeepSeekMath-RL 7B, GSM8K and MATH with chain-of-thought reasoning are treated as in-domain tasks, whereas all other benchmarks are considered out-of-domain tasks.

Table \ref{tab:sft_rl_math} presents the results of open- and closed-source models evaluated using both chain-of-thought and tool-assisted reasoning on English and Chinese benchmark datasets. Our observations are as follows: 1) DeepSeekMath-RL 7B achieves accuracies of 88.2\% on GSM8K and 51.7\% on MATH when employing chain-of-thought reasoning. This level of performance exceeds that of all open-source models ranging from 7B to 70B parameters, as well as most closed-source models. 2) Importantly, DeepSeekMath-RL 7B is trained solely on chain-of-thought formatted instruction tuning data from GSM8K and MATH, initialized from DeepSeekMath-Instruct 7B. Despite the limited scope of its training dataset, it surpasses DeepSeekMath-Instruct 7B in every evaluation metric, demonstrating the advantage of reinforcement learning.

\section{Discussion}
In this section, we discuss the insights gained from our pre-training and reinforcement learning experiments.

\subsection{Insights Gained from Pre-Training}

We begin by sharing our experiences related to pre-training. Unless stated otherwise, the training configurations follow those described in Section \ref{sec:quality-policy}. Note that when we mention the DeepSeekMath Corpus here, we refer to the 89-billion-token dataset collected during the second iteration of data acquisition.

\subsubsection{Training on Code Enhances Mathematical Reasoning}

A widely held yet unconfirmed belief is that training on code improves reasoning abilities. We aim to provide partial evidence supporting this notion specifically in the domain of mathematics: incorporating code training enhances the model’s capacity for mathematical reasoning, both when using tools and independently.

To investigate the influence of code training on mathematical reasoning, we conducted experiments using two training approaches: two-stage training and one-stage training:

\textbf{Two-Stage Training}

\begin{itemize}[topsep=0pt]
    \item \textbf{Code Training for 400B Tokens $\rightarrow$ Math Training for 150B Tokens}: We first train DeepSeek-LLM 1.3B on 400B code tokens, followed by an additional 150B math tokens;
    \item \textbf{General Training for 400B Tokens $\rightarrow$ Math Training for 150B Tokens}: As a baseline, we also conduct training using general tokens (drawn from a large-scale general corpus curated by DeepSeek-AI) instead of code tokens in the initial training phase, aiming to explore whether code tokens offer advantages over general tokens in enhancing mathematical reasoning.
\end{itemize}

\textbf{One-Stage Training}

\begin{itemize}[topsep=0pt]
    \item \textbf{Math Training for 150B Tokens}: DeepSeek-LLM 1.3B is trained solely on 150B math tokens;
    \item \textbf{Training on a mixture of 400B Code Tokens and 150B Math Tokens}: Since conducting math training after code training reduces coding performance, we examine if blending code and math tokens in a single-stage training regimen can still improve mathematical reasoning while also reducing the problem of catastrophic forgetting.
\end{itemize}

\paragraph{Results}
Tables \ref{tab:code-math} and \ref{tab:code-math-general-eval} present downstream performance outcomes across various training configurations.

Code training enhances program-assisted mathematical reasoning in both two-stage and one-stage training approaches. As indicated in Table \ref{tab:code-math}, for the two-stage method, training exclusively on code tokens substantially improves the model’s capability to solve GSM8K and MATH problems using Python. Following up with math training in the second stage further boosts performance. Notably, in the one-stage training scenario, mixing code and math tokens effectively addresses the catastrophic forgetting problem seen in two-stage training, while also promoting synergy between coding (Table \ref{tab:code-math-general-eval}) and program-aided mathematical reasoning (Table \ref{tab:code-math}).

Code training also advances mathematical reasoning without reliance on tools. Within the two-stage training setup, the initial code training phase already produces moderate gains, which enhance the efficacy of the later math training, ultimately achieving the highest performance. Conversely, when code and math tokens are combined in a single training stage, mathematical reasoning without tools deteriorates. One possible explanation is that the 1.3B parameter size of DeepSeek-LLM limits its ability to effectively integrate both code and mathematical content simultaneously.

\subsubsection{ArXiv Papers Appear Ineffective in Enhancing Mathematical Reasoning}

ArXiv papers are frequently used as a part of mathematical pre-training datasets \citep{minerva,gpt-f,llemma,mathpile}. Nonetheless, an in-depth evaluation of their influence on mathematical reasoning has been limited. Surprisingly, based on our experimental results, arXiv papers do not seem to significantly enhance mathematical reasoning capabilities. We tested models of various sizes, including DeepSeek-LLM 1.3B and DeepSeek-Coder-Base-v1.5 7B \citep{deepseek-coder}, utilizing arXiv corpora that were processed through different pipelines:

\begin{itemize}[topsep=0pt]
    \item \textbf{MathPile} \citep{mathpile}: an 8.9B-token dataset curated with cleaning and filtering heuristics, where more than 85\% consists of scientific arXiv papers;
    \item \textbf{ArXiv-RedPajama} \citep{redpajama}: the complete set of arXiv LaTeX files with preambles, comments, macros, and bibliographies removed, totaling 28.0B tokens.
\end{itemize}

In our trials, we independently trained DeepSeek-LLM 1.3B for 150B tokens and DeepSeek-Coder-Base-v1.5 7B for 40B tokens on each arXiv corpus. The findings indicate that arXiv papers do not effectively improve mathematical reasoning. Models trained solely on arXiv-based corpora exhibited either negligible gains or even declines in performance across a variety of mathematical benchmarks of varying difficulty utilized in this work. These benchmarks include quantitative reasoning datasets such as GSM8K and MATH (Table \ref{tab:arxiv-cot}), multiple-choice exams like MMLU-STEM (Table \ref{tab:arxiv-cot}), and formal mathematics tasks including miniF2F (Table \ref{tab:arxiv_atp}).

However, this conclusion comes with caveats and should be interpreted cautiously. We have yet to investigate:

\begin{itemize}[topsep=0pt]
    \item The influence of arXiv tokens on specific math-related tasks outside the scope of this study, such as theorem informalization, which involves converting formal statements or proofs into informal expressions;
    \item The effects when arXiv tokens are combined with other data types;
    \item Whether advantages from arXiv papers emerge at larger model scales.
\end{itemize}

Hence, additional research is necessary, which we reserve for future work.

\subsection{Insights of Reinforcement Learning}
\subsubsection{Towards a Unified Paradigm} In this section, we propose a unified framework to analyze various training strategies, including SFT, RFT, DPO, PPO, GRPO, and conduct experiments to investigate elements of this unified paradigm. Generally, the gradient with respect to the parameter $\theta$ for a training method can be expressed as:

\begin{equation}
    \nabla_{\theta}\mathcal{J}_{\textcolor{red}{\mathcal{A}}}(\theta) = \mathbb{E}[\underbrace{(q,o) \sim \textcolor{red}{\mathcal{D}}}_{Data \ Source}]\left( \frac{1}{|o|} \sum_{t=1}^{|o|}  \underbrace{GC_{{\mathcal{A}}}(q, o, t, \textcolor{red}{\pi_{{rf}}})}_{Gradient \ Coefficient}  \nabla_{\theta}\log \pi_{\theta}(o_t | q, o_{<t})\right).
\label{eq:objective}
\end{equation}

This framework includes three fundamental elements: 1) \textit{Data Source $\mathcal{D}$}, which specifies the dataset used for training; 2) \textit{Reward Function $\pi_{{rf}}$}, which provides the reward signal for learning; 3) \textit{Algorithm $\mathcal{A}$}: which transforms the training data and reward signal into the gradient coefficient $GC$, controlling the strength of the penalty or reinforcement applied to the data. We examine various representative methods within this unified framework:

\begin{itemize}
    \item  \textbf{Supervised Fine-tuning (SFT)}: SFT adapts a pretrained model using human-curated SFT datasets.
    \item \textbf{Rejection Sampling Fine-tuning (RFT)}: RFT further adjusts the SFT model by fine-tuning it on outputs filtered from samples generated by the SFT model based on SFT queries. The filtering is done by assessing the correctness of the answers.
    \item \textbf{Direct Preference Optimization (DPO)}: DPO refines the SFT model by fine-tuning on augmented outputs from the SFT model, utilizing a pairwise DPO loss function.
    \item \textbf{Online Rejection Sampling Fine-tuning (Online RFT)}: Unlike RFT, Online RFT begins with the SFT model as the policy and improves it by fine-tuning on augmented outputs sampled from the current policy model in real time.
    \item \textbf{PPO/GRPO}: PPO/GRPO starts with the SFT model to initialize the policy and enhances it by reinforcing outputs sampled from the live policy model.
\end{itemize}

We provide a summary of the components of these methods in Table \ref{tab:data_policy}. For a more comprehensive derivation process, please consult Appendix \ref{app:analysis-rl}.

\paragraph{Observation about Data Source} We categorize the data sources into two types: online sampling and offline sampling. Online sampling refers to training data generated from exploration outcomes of the current training policy model, whereas offline sampling refers to data obtained from the sampling results of the initial SFT model. RFT and DPO utilize the offline approach, while Online RFT and GRPO adopt the online approach.

As illustrated in Figure \ref{fig:rl-analysis}, we observe that Online RFT considerably outperforms RFT across two benchmarks. Notably, Online RFT performs similarly to RFT during the early training phases but gains a clear advantage in later stages, demonstrating the benefits of online training. This observation is intuitive because, initially, the actor and the SFT model are quite similar, and the sampled data shows only minor variations. However, as training progresses, data sampled from the actor increasingly diverges, making real-time data sampling more beneficial.

\paragraph{Observation about Gradient Coefficient} The algorithm converts the reward signal into a gradient coefficient to update the model parameters. In our experiments, we classify the reward function into two types: `Rule' and `Model'. The `Rule' category involves evaluating the response quality based on answer correctness, whereas the `Model' category involves training a reward model to assign scores to each response. The reward model is trained using data labeled according to the rule-based judgment. Equations \ref{eq:GC-RFT} and \ref{eq:GC-GRPO} emphasize a critical distinction between GRPO and Online RFT: GRPO uniquely modulates its gradient coefficient according to the reward value from the reward model, allowing it to differentially reinforce or penalize responses based on varying magnitudes. By contrast, Online RFT does not implement this feature; it does not penalize incorrect responses and applies uniform reinforcement to all correct responses at the same intensity level.

As illustrated in Figure \ref{fig:rl-analysis}, GRPO outperforms online RFT, underscoring the effectiveness of adjusting positive and negative gradient coefficients. Moreover, GRPO+PS demonstrates better results than GRPO+OS, suggesting the advantages of employing fine-grained, step-aware gradient coefficients. Additionally, we investigate iterative RL by conducting two rounds of iteration in our experiments. As depicted in Figure \ref{fig:iter-rl}, iterative RL substantially enhances performance, particularly after the first iteration.

\subsubsection{Why Does RL Work?} In this study, reinforcement learning is applied to a subset of instruction tuning data, leading to a notable improvement over the instruction tuning model. To further clarify why reinforcement learning is effective, we assess the Pass@K and Maj@K accuracy metrics of the Instruct and RL models on two benchmarks. As presented in Figure \ref{fig:maj-pass}, RL boosts Maj@K performance but not Pass@K. This suggests that RL improves the model’s overall output robustness by increasing the likelihood of the correct answer appearing among the TopK outputs, rather than enhancing the model’s core capabilities. Correspondingly, \citep{wang2023making} identified a \textbf{misalignment issue} in reasoning tasks within SFT models, showing that reasoning can be improved through a series of preference alignment methods \citep{yuan2023rrhf,song2023preference,wang2023making}.

\subsubsection{How Can RL Be Made More Effective?}
\label{sec:effec_rl}
We show that RL is quite effective in mathematical reasoning tasks. Additionally, we propose a unified framework to interpret various representative training methods, viewing all of them as either direct or simplified RL approaches. As outlined in Equation \ref{eq:objective}, three main components are involved: Data Source, Algorithm, and Reward Function. We suggest several possible future research directions related to these components.

\paragraph{Data Source} The data source serves as the foundation for all training methods. For RL, this specifically refers to unlabeled questions with outputs generated by the policy model. In this work, we utilize only questions from the instruction tuning phase and apply a straightforward nucleus sampling technique to produce outputs. We consider this a possible reason why our RL framework mainly enhances Maj@K performance. Going forward, we plan to test our RL pipeline on out-of-distribution question prompts combined with \textbf{advanced sampling (decoding) strategies}, such as those based on tree-search methods \citep{yao2023tree}. Furthermore, \textbf{efficient inference techniques} \citep{xia-etal-2023-speculative,leviathan2023fast,kwon2023efficient,xia2024unlocking}, which govern the exploration efficiency of policy models, will also play a crucial role.

\paragraph{Algorithms} Algorithms utilize the data and reward signals to compute the gradient coefficient, which in turn updates the model parameters. According to Equation \ref{eq:objective}, currently all methods to some degree fully \textbf{TRUST} the reward function's signal to adjust the conditional probability of specific tokens. Nonetheless, it is unrealistic to guarantee that the reward signal is always dependable, particularly in highly complex tasks. For instance, even the PRM800K dataset \citep{lightman2023let}, meticulously annotated by expert annotators, contains about 20\% incorrect annotations\footnote{\url{https://github.com/openai/prm800k/issues/12\#issuecomment-1728491852}}. In light of this, we aim to investigate reinforcement learning algorithms that remain robust in the presence of noisy reward signals. We posit that such \textbf{WEAK-TO-STRONG} \citep{burns2023weak} alignment approaches could fundamentally transform learning algorithms.

\paragraph{Reward Function} The reward function serves as the foundation for the training signal. In reinforcement learning, the reward function is typically realized by a neural reward model. We identify three critical avenues for advancing reward models: 1) \textbf{Enhancing the generalization capacity of the reward model.} It is essential for the reward model to generalize effectively to out-of-distribution queries and sophisticated decoding outputs; otherwise, reinforcement learning might only maintain the existing distribution of LLMs instead of fundamentally improving their capabilities; 2) \textbf{Representing the uncertainty inherent in the reward model.} This uncertainty could potentially act as a crucial connection between weak reward models and weak-to-strong learning algorithms; 3) \textbf{Developing high-quality process reward models efficiently} that deliver detailed training signals specifically for the reasoning process \citep{lightman2023let,wang2023math}.

\section{Conclusion, Limitation, and Future Work} We introduce DeepSeekMath, which surpasses all open-source models on the competition-level MATH benchmark and nears the performance of closed-source models. DeepSeekMath~is initialized from DeepSeek-Coder-v1.5 7B and is continually trained on 500B tokens, with a substantial portion of the training data being 120B math tokens obtained from Common Crawl. Our comprehensive ablation study reveals that web pages present significant opportunities for acquiring high-quality mathematical data, whereas arXiv may not be as advantageous as initially anticipated. We propose Group Relative Policy Optimization (GRPO), a variant of Proximal Policy Optimization (PPO), which significantly enhances mathematical reasoning abilities while reducing memory usage. Experimental results demonstrate that GRPO remains effective even though DeepSeekMath-Instruct 7B has already achieved high benchmark scores. Additionally, we offer a unified framework to interpret a range of methods and outline several promising directions for more efficient reinforcement learning.

Although DeepSeekMath~achieves strong results on quantitative reasoning benchmarks, its performance on geometry and theorem-proving tasks is comparatively weaker than that of closed models. For example, during our preliminary testing, the model struggled with problems involving triangles and ellipses, which may suggest a bias in data selection during pre-training and fine-tuning. Moreover, constrained by model size, DeepSeekMath~underperforms compared to GPT-4 in few-shot capabilities. GPT-4 improves with few-shot input, whereas DeepSeekMath~exhibits similar performance in both zero-shot and few-shot evaluations. Moving forward, we plan to refine our engineered data selection pipeline to build a higher-quality pre-training corpus. Furthermore, we intend to investigate the potential directions outlined in Section \ref{sec:effec_rl} to develop more effective reinforcement learning strategies for large language models.

\end{document}